\chapter{\textcolor{mycolor}{\textbf{Tests, validation et analyse des performances}}}
\markboth{\chaptername\ \thechapter.\ Tests, validation et analyse des performances}{}
\fancyhf{}
\fancyhead[L]{\small\itshape\nouppercase{\rightmark}}
\fancyfoot[L]{\small\itshape\nouppercase{\leftmark}}
\fancyfoot[R]{\small\itshape\thepage}
\renewcommand{\headrulewidth}{1pt}
\renewcommand{\footrulewidth}{0.4pt}

\vfill
\begin{flushright}{\Large\bfseries\textcolor{mycolor}{Introduction}}\end{flushright}
\addcontentsline{toc}{section}{Introduction}

Ce chapitre rend compte de l'évaluation du prototype. Nous décrivons le protocole expérimental, puis les résultats obtenus pour le matériel, la vision, la décision et l'usage du tableau de bord. Une discussion générale rapproche ces résultats des objectifs fixés au premier chapitre.

\clearpage

% Sous-équipe 2 : Data. Section du chapitre 5, appelée par rapport/main chapters.
% Retirer l'encadré \aRediger une fois la section écrite.

\section{Protocole expérimental}

\aRediger{2 (Data)}{Décrire les scénarios de test, la vérité terrain et les métriques retenues.}

\subsection{Scénarios de test}

\subsection{Constitution de la vérité terrain}

\subsection{Choix des métriques}


% Sous-équipe 1 : IoT et matériel. Section du chapitre 5, appelée par rapport/main chapters.
% Retirer l'encadré \aRediger une fois la section écrite.

\section{Performances matérielles}

\aRediger{1 (IoT et matériel)}{Présenter les mesures matérielles relevées pendant la démonstration.}

\subsection{Temps d'acquisition}

\subsection{Processeur et mémoire sur le Raspberry Pi}


% Sous-équipe 3 : Vision (VLM). Section du chapitre 5, appelée par rapport/main chapters.
% Retirer l'encadré \aRediger une fois la section écrite.

\section{Évaluation du module de vision}

\aRediger{3 (Vision (VLM))}{Présenter les résultats (objets bien vus, manqués, inventés), l'écart entre les deux montages et l'analyse des erreurs.}

\subsection{Résultats par modèle et par caméra}

\subsection{Comparaison ordinateur et Raspberry Pi}

\subsection{Analyse des erreurs}


% Sous-équipe 4 : Décision et backend. Section du chapitre 5, appelée par rapport/main chapters.
% Retirer l'encadré \aRediger une fois la section écrite.

\section{Évaluation des décisions}

\aRediger{4 (Décision et backend)}{Présenter le taux de décisions correctes, les fausses alertes et la latence mesurée.}

\subsection{Décisions correctes et fausses alertes}

\subsection{Latence de bout en bout}


% Sous-équipe 5 : Application web. Section du chapitre 5, appelée par rapport/main chapters.
% Retirer l'encadré \aRediger une fois la section écrite.

\section{Tests d'usage}

\aRediger{5 (Application web)}{Présenter les tests d'usage du tableau de bord : historique, suppression, réglages.}

\subsection{Protocole des tests d'usage}

\subsection{Résultats}


% Sous-équipe 6 : Application mobile et coordination. Section du chapitre 5, appelée par rapport/main chapters.
% Retirer l'encadré \aRediger une fois la section écrite.

\section{Discussion générale}

\aRediger{6 (Application mobile et coordination)}{Rapprocher les résultats des objectifs du chapitre 1, présenter les difficultés et les limites.}

\subsection{Synthèse des résultats}

\subsection{Comparaison avec les objectifs}

\subsection{Difficultés rencontrées}

\subsection{Limites de l'expérimentation}



\section*{Conclusion}
\addcontentsline{toc}{section}{Conclusion}

\aRediger{6 (coordination)}{Conclure le chapitre en une dizaine de lignes.}
